\section{Installation and User Manual of the Word-Meaning Lab}
\label{app:manual}

\paragraph{Installation.}
The reference environment is Python 3.12 with the package versions pinned in \texttt{requirements.txt} of the project
repository, including Streamlit.
From the repository root:
\begin{quote}\small\ttfamily
python3.12 -m venv .venv\\
.venv/bin/python -m pip install -r requirements.txt\\
.venv/bin/streamlit run app/word\_lab.py
\end{quote}
Streamlit prints a local address and opens it in the browser.
The application needs only the committed folder \texttt{results/\allowbreak words/}; if a required file is missing,
the page names it together with the script that produces it.
The logic tests and the smoke tests of both pages run with
\begin{quote}\small\ttfamily
.venv/bin/python -I -m pytest -q tests/test\_word\_lab.py
\end{quote}
Newer Python versions may need the workaround documented in \texttt{app/README.md}.

\paragraph{First page: What does a word mean?}
\begin{enumerate}
\item Choose a phrase in the drop-down list.
\item Read the chart: the horizontal axis is the probability in percent, the height is the membership, that is, how well
each number fits the phrase. Blue outlines show the citizen codebook: the outer line is the widest demographic group,
the inner line is where all groups agree, and the shading between them shows disagreement. Orange outlines show the
LLM codebook in the same way. Dashed lines show each model; a tall thin spike means a model always gave the same number.
\item Use the switches to show or hide the model lines and the six group lines.
\item Read the headline numbers: overlap (Jaccard, from zero to one), centre shift (LLM minus respondents, in
percentage points) and the number of models that gave one fixed number; hover over the help icon for the confidence
interval. A note appears when a model was excluded for the phrase.
\item The tab ``All words ranked'' orders all phrases by overlap, with intervals, and marks hedge and extreme words.
\end{enumerate}

\paragraph{Second page: Crisis-line simulator.}
\begin{enumerate}
\item Set the crisis degree, from zero to one hundred: in the study, the share of LLM labellers that marked a message
as a crisis.
\item Choose the quality of the bot's reply, from \word{Harmful/Unacceptable} to \word{Fully Appropriate}.
\item Read the banner: red ``Decision flips'' if the citizen and the LLM codebook lead to different actions, with the
direction; green ``Same decision'' otherwise.
\item The table gives, for each codebook, the crisis word, its fit and the action (A0 continue, A1 suggest a helpline,
A2 hand over to a counsellor, A3 urgent handover with emergency numbers); highlighted rows differ from the citizen
codebook, and ``no fit'' marks degrees that no word covers.
\item The expander ``Where does each crisis word start and stop?'' shows the four crisis words of both codebooks with the
slider position.
\end{enumerate}

\paragraph{Known limitations.}
The application is a research prototype and not a clinical tool; it shows only prompt P1, only the four models of the
study and only the six demographic groups; it shows the study's confidence intervals but does not recompute them; and
it has not been tested with a screen reader or with users from crisis services.
